\section{Legal Requirements for Congressional Redistricting}
\label{sec:legal-requirements}

Congressional redistricting must satisfy a hierarchy of legal requirements
derived from the federal Constitution, the Voting Rights Act, and state
constitutions and statutes.
We enumerate eleven requirements that any redistricting algorithm must address,
organising them from the hardest federal floor to the softer state-specific ceilings.
Each requirement has a precise algorithmic interpretation that Section~\ref{sec:requirements-matrix}
will map to specific toolbox components.

\subsection{Constitutional Floor Requirements}

\subsubsection*{R1: Population Equality ($\pm 0.5\,\%$)}

\textbf{Legal basis}: \wesberry{} v.\ Sanders, 376 U.S.\ 1 (1964) (``one person,
one vote''); Karcher v.\ Daggett, 462 U.S.\ 725 (1983) (near-perfect equality
required for congressional districts).

Congressional districts must contain as close to equal population as mathematically
possible.
Unlike state legislative districts, which allow deviations up to $\pm 10\,\%$
under \emph{Mahan v.\ Howell} (1973), congressional districts face a near-zero
tolerance: any deviation must be ``unavoidable'' or justified by a legitimate
state policy.

\textbf{Algorithmic interpretation}: All METIS calls use $\mathtt{ubvec}[0] = 1.001$
(0.1\,\% tolerance per part, yielding $\leq 0.5\,\%$ maximum deviation across $k$
parts by the triangle inequality).
Population balance is the hard constraint; all other objectives are soft.

\subsubsection*{R2: Contiguity}

\textbf{Legal basis}: Article~I, \S2 (no explicit text, but universally treated
as a constitutional floor by courts); all 50 state constitutions require contiguous
districts.

Every district must be a connected geographic region.
No district may consist of non-adjacent parcels even if they share the same
population.

\textbf{Algorithmic interpretation}: METIS's \texttt{-contig} flag is set at
every bisection level.
For planar census-tract adjacency graphs, METIS empirically produces contiguous
partitions even without this flag \citep{b8paper}, but the flag is set explicitly
to provide a formal guarantee.

\subsection{Compactness Requirements}

\subsubsection*{R3: Compactness}

\textbf{Legal basis}: Most state constitutions (California Constitution Art.\ XXI
\S2(d); New York ELEC \S4-100; Iowa Code \S42.4); some states treat compactness
as a federal equal-protection requirement under \shaw{} v.\ Reno, 509 U.S.\ 630
(1993) (``bizarre'' shapes trigger heightened scrutiny).

Districts should be reasonably compact in shape: not elongated, not fractal,
not reaching tentacles into remote areas to capture or avoid specific communities.
Courts have measured compactness via the Polsby-Popper score (ratio of district
area to area of circumscribed circle), the Reock score (ratio to minimum
bounding circle), and the isoperimetric ratio (perimeter squared over area).

\textbf{Algorithmic interpretation}: Geographic edge weights (B.2) encode TIGER
boundary lengths as edge costs, making it expensive for METIS to cut across long
geographic boundaries.
The isoperimetric normalisation of GeoSection (B.8) additionally prevents the
caterpillar pathology (sequential urban peeling) that creates compactness violations
invisible to per-district metrics but visible in the bisection tree structure.

\subsection{Subdivision Preservation}

\subsubsection*{R4: County and Subdivision Preservation}

\textbf{Legal basis}: Rucho v.\ Common Cause, 588 U.S.\ 684 (2019) (partisan
gerrymandering nonjusticiable, but traditional criteria---including county
preservation---remain valid state-court grounds); Harper v.\ Hall (N.C. Sup.\ Ct.,
2022) (required county grouping); most state redistricting statutes list
``preserving political subdivisions'' as an explicit criterion.

Districts should avoid splitting counties, municipalities, and townships where
possible, consistent with the population-equality requirement.
County preservation has two rationales: it minimises disruption to existing
community administrative structures, and it provides a measurable baseline that
prevents ``county-splitting gerrymandering'' where a partisan actor systematically
splits unfriendly counties to dilute their political influence.

\textbf{Algorithmic interpretation}: County-sticky edge weights (B.10) multiply
edge costs by $(1 + \alpha_\text{county})$ for pairs of census tracts sharing the
same county FIPS code.
As $\alpha_\text{county} \to \infty$, county splits occur only when the
population-equality constraint forces them.
The $\alpha$ parameter is a published number; a court can read it from the
plan manifest.

\subsection{Voting Rights Act Requirements}

\subsubsection*{R5: VRA \S2 Gingles Prong 1 --- Minority Geographic Compactness}

\textbf{Legal basis}: 52 U.S.C.\ \S10301; Thornburg v.\ Gingles, 478 U.S.\ 30
(1986) (three-prong test: minority group sufficiently large and geographically
compact; minority cohesion; majority bloc voting); Allen v.\ Milligan,
598 U.S.\ 1 (2023) (Alabama required two Black-opportunity districts);
\callais{}, 608 U.S.\ \_\_\_ (2026) (Gingles Prong 1 is geographic, not racial
targeting).

\S2 requires states to draw majority-minority districts where a minority
community is (1) sufficiently large and geographically compact to form a majority
in a single district, (2) politically cohesive, and (3) faces bloc voting by
the majority.
The geographic concentration prong (Prong 1) is the one that algorithmic
redistricting can directly address: the algorithm must be capable of identifying
whether a compact minority community exists at the geographic scale of a
congressional district.

\textbf{Algorithmic interpretation}: VRASection (B.14) adds a geographic alignment
score $A = |\text{MVAP}_\text{frac}(L) - \text{MVAP}_\text{frac}(R)|$ to the
ratio selection objective.
This score is a \emph{geographic property}: it measures whether the bisection
places minority population on one side.
Critically, the score uses only the spatial distribution of minority VAP,
not racial identity or any racial target.
At $w_\text{vra} = 0.40$, compactness is weighted 50\,\% more heavily than
alignment, ensuring racial geometry does not predominate.

\subsubsection*{R6: No Racial Predominance}

\textbf{Legal basis}: \shaw{} v.\ Reno, 509 U.S.\ 630 (1993); \miller{} v.\
Johnson, 515 U.S.\ 900 (1995) (race cannot be the ``predominant factor'' in
drawing district lines); \callais{} p.~36 (race-conscious and partisan signals
must be disentangled).

A district whose shape is ``so bizarre'' that it can only be explained by race
triggers strict scrutiny under \shaw{}.
Under \miller{}, any time race is the predominant factor overriding traditional
criteria (compactness, contiguity, political subdivisions), the district
requires a compelling interest.

\textbf{Algorithmic interpretation}: VRASection's $w_\text{vra} = 0.40$ weight
is calibrated to fall well below the \shaw{}--\miller{} racial-predominance
threshold.
The score does not specify a racial target (no ``district must be 50\,\% Black VAP'');
it uses only the geographic alignment of existing population distribution.
The design ensures that compactness remains the predominant objective
(see Section~\ref{sec:callais} for the formal analysis).

\subsection{Partisan Neutrality}

\subsubsection*{R7: Partisan Neutrality (State Courts)}

\textbf{Legal basis}: \lwvpa{} (Pa.\ Sup.\ Ct., 2018) (partisan gerrymander
violates Pennsylvania Constitution); Harper v.\ Hall (N.C.\ Sup.\ Ct., 2022)
(initial decision found unconstitutional partisan gerrymander); Harkenrider v.\
Hochul (N.Y.\ Ct.\ App., 2022) (struck down partisan gerrymander under
New York Constitution).

While \rucho{} (2019) closed the federal door to partisan gerrymandering claims,
state courts in Pennsylvania, North Carolina, New York, and other states have
entertained partisan-neutrality claims under state constitutional grounds.
The legal standard varies: some courts require proportionality of seats to votes,
others require map-making processes that do not give undue advantage.

\textbf{Algorithmic interpretation}: The Rodden-null test is the key tool here.
The B.9 AreaSection paper demonstrated that AreaSection and GeoSection produce
identical seat counts in 76\,\% of states and differ by only 1 seat in the
remaining states, with no systematic partisan direction.
The partisan outcomes in all competitive states are explained by geographic
voter sorting (the ``Rodden effect'' \citep{rodden2019}), not by algorithmic bias.
ProportionalSection (B.12, planned) addresses this directly by incorporating
partisan vote share as a separate, Callais-compliant mode.

\subsection{Cross-Census Stability}

\subsubsection*{R8: Census Stability (Litigation Argument)}

\textbf{Legal basis}: Not a statutory requirement, but a powerful litigation
argument against partisan gerrymandering claims.
If a redistricting algorithm produces the same map---the same natural bisection
ratio, the same seat counts---across 2000, 2010, and 2020 Census data, then
the map is not a product of the current census cycle's partisan distribution.
A challenger cannot argue the algorithm exploited the current population counts
to achieve partisan ends when the identical map would have been drawn from
twenty-year-old data.

\textbf{Algorithmic interpretation}: StabilitySection (B.15) measures the Census
Stability Score (CSS): a composite of seat-count stability, ratio stability, and
proportionality-gap stability across census years.
The 2000--2010 sweep found 67\,\% of states are stable at the $\Delta f \leq 0.05$
threshold; Iowa is the canonical unstable state.
A CSS $\geq 0.67$ threshold is proposed as the minimum for making the
census-stability legal argument.

\subsection{Auditability and Seed Invariance}

\subsubsection*{R9: Seed and Auditability Invariance}

\textbf{Legal basis}: No explicit legal requirement, but a pervasive requirement
in redistricting litigation: courts need to be able to reproduce the map from
the stated parameters.
A plan produced by an algorithm whose outputs vary with undisclosed random seeds
is not auditable; a challenger can always claim the seed was chosen to achieve
a specific partisan outcome.

\textbf{Algorithmic interpretation}: The B.7 paper certified that GeoSection
achieves near-zero partisan variance across random seeds (seat-count variance
$< 0.3$ seats per state at 50 seeds per ratio).
ApportionRegions (B.11, planned) extends this to the Huntington-Hill linkage:
the natural bisection ratio can be derived from the state's prime factorisation
of district counts, making it parameter-free and auditable.
The \texttt{redist} CLI records every parameter in the \texttt{whatif-manifest v1}
JSON sidecar, and a hash chain verifies that no parameter was altered after
the run.

\subsection{Multi-Chamber Consistency}

\subsubsection*{R10: Multi-Chamber Consistency}

\textbf{Legal basis}: State constitutions require that congressional and state
legislative district maps be drawn consistently; the ``one-map'' approach prevents
``stacked'' gerrymandering where a partisan actor controls all three chambers and
packs or cracks a community at every level.

\textbf{Algorithmic interpretation}: NestSection (B.13) forces all chambers
sharing a state to use a \emph{compatible factorisation spine} derived from
$g = \gcd(C, S, H)$ where $C$, $S$, $H$ are the congressional, senate, and
house seat counts.
Eleven states achieve exact nesting (spine compatibility score 0); the remainder
require best-effort Mode~3 nesting.
A shared spine is an anti-gerrymandering mechanism: a partisan actor who locks
in a congressional tree cannot then draw senate boundaries that contradict it.

\subsection{Huntington-Hill Linkage}

\subsubsection*{R11: Apportionment Linkage}

\textbf{Legal basis}: Article~I, \S2 (apportionment proportional to population);
the Huntington-Hill method (2 U.S.C.\ \S2a) as the statutory apportionment
algorithm.
A redistricting algorithm that uses a bisection tree structurally linked to the
state's Huntington-Hill seat count provides an additional layer of constitutional
grounding: the geographic partition reflects the same proportionality principle
as the apportionment itself.

\textbf{Algorithmic interpretation}: ApportionRegions (B.11, planned) determines
the bisection tree structure from the prime factorisation of $k$, mirroring the
successive subdivision implicit in the Huntington-Hill priority sequence.
A state with $k = 8 = 2^3$ can be exactly bisected three levels deep.
A state with $k = 7$ (prime) requires a near-equal first split and defers the
odd seat to the final level.
The linkage is not a hard constitutional requirement but provides additional
grounding for the algorithm's objectivity claim in litigation.
